\documentclass{article}
\usepackage{fontspec}
\setmainfont{Libertinus Serif}
\usepackage[a4paper, margin=1in]{geometry}
\usepackage{amsmath, amsthm, amssymb, mathtools, enumitem, hyperref}

\theoremstyle{plain}
\newtheorem{thm}{Theorem}[section]
\newtheorem{prop}[thm]{Proposition}
\newtheorem{lem}[thm]{Lemma}
\newtheorem{cor}[thm]{Corollary}

\theoremstyle{definition}
\newtheorem{defn}{Definition}[section]
\newtheorem{exmp}[defn]{Example}
\newtheorem{prblm}[defn]{Problem}
\newtheorem{prcp}[defn]{Principle}

\theoremstyle{remark}
\newtheorem*{rem}{Remark}

\DeclareMathOperator{\spn}{span}

\geometry{
 a4paper,
 left=3cm,
 top=2cm,
 right=2cm,
 bottom=2cm
}

\begin{document}
\begin{center}
    \textbf{\large UNIVERSIDADE FEDERAL DO RIO GRANDE DO SUL} \\
    \textbf{\large INSTITUTO DE MATEMÁTICA E ESTATÍSTICA} \\
    \vspace{3cm}
    \textbf{\Large ÁLGEBRA LINEAR I} \\
    \vspace{1cm}
    \textbf{\large LISTA DE DEFINIÇÕES E RESULTADOS} \\
    \vspace{3cm}

    \textbf{LEONARDO SUAREZ}

    \vfill
    \textbf{\large Porto Alegre, RS} \\
    \textbf{\large \today}
\end{center}

\newpage
\tableofcontents
\newpage

\noindent From the sixth edition of \textit{Linear Algebra and Its Applications}, by David Lay, Steven R. Lay, and Judi J. McDonald.

\section{Linear Equations in Linear Algebra}

\begin{thm} (Uniqueness of the Reduced Echelon Form). 
    Each matrix is row equivalent to one and only one reduced echelon matrix.
\end{thm}

\begin{thm} (Uniqueness and Existence Theorem).
    A linear system is consistent iff the rightmost column of the augmented matrix is \textit{not} a pivot column---that is, iff an echelon form of the augmented matrix has \textit{no} row of the form \[[0\hspace{.25cm}\cdots\hspace{.25cm}0\hspace{.25cm} b],\hspace{.5cm}b\neq0\] If a linear system is consistent, then the solution set contains either (i) a unique solution, when there are no free variables, or (ii) infinitely many solutions, when there is at least one free variable.
\end{thm}

\begin{thm}
    If $A$ is an \(m\times n\) matrix, with columns $\textbf{a}_1,\ldots,\textbf{a}_n$, and if $\textbf{b}$ is in $\mathbb{R}^m$, the matrix equation $$A\textbf{x}=\textbf{b}$$ has the same solution set as the vector equation $$x_1\textbf{a}_1+x_2\textbf{a}_2+\cdots+x_n\textbf{a}_n$$ which, in turn, has the same solution set as the system of linear equations whose augmented matrix is
    \[
    \begin{bmatrix}
        \textbf{a}_1 & \textbf{a}_2 & \cdots & \textbf{a}_n & \textbf{b}
    \end{bmatrix}
    \]
\end{thm}

\begin{thm}
    Let $A$ be an $m\times n$ matrix. Then the following statements are logically equivalent.
    \begin{enumerate}
        \item[a.] For each $\textbf{b}$ in $\mathbb{R}^m$, the equation $A\textbf{x}=\textbf{b}$ has a solution;
        \item[b.] Each $\textbf{b}$ is a linear combination of the columns of $A$;
        \item[c.] The columns of $A$ span $\mathbb{R}^m$;
        \item[d.] $A$ has a pivot in every row.
    \end{enumerate}
\end{thm}

\begin{thm}
    If $A$ is an $m\times n$ matrix, $\textbf{u}$ and $\textbf{v}$ are vectors in $\mathbb{R}^m$, and $c$ is a scalar, then:
    \begin{enumerate}
        \item[a.] $A(\textbf{u}+\textbf{v})=A\textbf{u}+A\textbf{v}$;
        \item[b.] $A(c\textbf{u})=c(A\textbf{u})$.
    \end{enumerate}
\end{thm}

\begin{thm}
    Suppose the equation $A\textbf{x}=\textbf{b}$ is consistent for some given $\textbf{b}$, and let $\textbf{p}$ be a solution. Then the solution set of $A\textbf{x}=\textbf{b}$ is the set of all vectors of the form $\textbf{w}=\textbf{p}+\textbf{v}_{h}$, where $\textbf{v}_{h}$ is any solution to the homogeneous equation $A\textbf{x}=\textbf{0}$.
\end{thm}

\begin{thm} (Characterization of Linearly Dependent Sets).
    An indexed set $S=\{\textbf{v}_1,\ldots,\textbf{v}_p\}$ of two or more vectors is linearly dependent iff at least one of the vectors in $S$ is a linear combination of the others. In fact, if $S$ is linearly dependent and $\textbf{v}_1\neq=\textbf{0}$, then some $\textbf{v}_j$ (with $j>1$) is a linear combination of the preceding vectors, $\textbf{v}_1,\ldots,\textbf{v}_{j-1}$.
\end{thm}

\begin{thm}
    Any set $\{\textbf{v}_1,\ldots,\textbf{v}_p\}$ in $\mathbb{R}^n$ is linearly dependent if $p>n$.
\end{thm}

\begin{thm}
    If a set $S=\{\textbf{v}_1,\ldots,\textbf{v}_p\}$ in $\mathbb{R}^n$ contains the zero vector, then the set is linearly dependent.
\end{thm}

\begin{thm}
    Let $T:\mathbb{R}^n\to\mathbb{R}^m$ be a linear transformation. Then there exists a unique matrix $A$ such that $$T(\textbf{x})=A\textbf{x},\hspace{.5cm}\forall\textbf{x}\in\mathbb{R}^n$$ In fact, $$A=[T(\textbf{e}_1)\hspace{.5cm}\cdots\hspace{.5cm}T(\textbf{e}_n)]$$    
\end{thm}

\begin{thm}
    Let $T:\mathbb{R}^n\to\mathbb{R}^m$ be a linear transformation. Then $T$ is one-to-one iff the equation $T(\textbf{x})=\textbf{0}$ has only the trivial solution.
\end{thm}

\begin{thm}
    Let $T:\mathbb{R}^n\to\mathbb{R}^m$ be a linear transformation, and let $A$ be the standard matrix for $T$. Then:
    \begin{enumerate}
        \item[a.] $T$ maps $\mathbb{R}^n$ onto $\mathbb{R}^m$ iff the columns of $A$ span $\mathbb{R}^m$;
        \item[b.] $T$ is one-to-one iff the columns of $A$ are linearly independent.
    \end{enumerate}
\end{thm}
    
\section{Matrix Algebra}

\begin{thm}
    Let $A$, $B$, and $C$ be matrices of the same size, and let $r$ and $s$ be scalars.
    \begin{enumerate}
        \item[a.] $A+B=B+A$
        \item[b.] $(A+B)+C=A+(B+C)$
        \item[c.] $A+0=A$
        \item[d.] $r(A+B)=rA+rB$
        \item[e.] $(r+s)A=rA+sA$
        \item[f.] $r(sA)=(rs)A$
    \end{enumerate}
\end{thm}

\begin{thm} (Properties of Matrix Multiplication).
    Let $A$ be an $m\times n$ matrix, and let $B$ and $C$ have sizes for which the indicated sums and products are defined.
    \begin{enumerate}
        \item[a.] $A(BC)=(AB)C$ \hfill(associative law of multiplication)
        \item[b.] $A(B+C)=AB+AC$ \hfill(left distributive law)
        \item[c.] $(B+C)A=BA+CA$ \hfill(right distributive law)
        \item[d.] $r(AB)=(rA)B=A(rB)$ \hfill for any scalar $r$
        \item[e.] $I_mA=A=AI_n$
    \end{enumerate}
\end{thm}

\begin{thm} (Properties of Matrix Transposes).
    Let $A$ and $B$ denote matrices whose sizes are appropriate for the following sums and products.
    \begin{enumerate}
        \item[a.] $(A^T)^T=A$
        \item[b.] $(A+B)^T=A^T+B^T$
        \item[c.] For any scalar $r$, $(rA)^T=rA^T$
        \item[d.] $(AB)^T=B^TA^T$
    \end{enumerate}
\end{thm}

\begin{thm}
    Let $A=\begin{bmatrix}
        a & b\\
        c & d
    \end{bmatrix}$. If $ad-bc\neq0$, then $A$ is invertible and $$A^{-1}=\frac{1}{ad-bc}\begin{bmatrix}
        d & -b\\
        -c & a
    \end{bmatrix}$$ If $ad-bc=0$, then $A$ is not invertible.
\end{thm}

\begin{thm}
    If $A$ is an invertible $n\times n$ matrix, then for each $\textbf{b}$ in $\mathbb{R}^n$, the equation $A\textbf{x}=\textbf{b}$ has the unique solution $\textbf{x}=A^{-1}\textbf{b}$.
\end{thm}

\begin{thm}
    \begin{enumerate}
        \item[a.] If $A$ is an invertible matrix, then $A^{-1}$ is invertible and $$(A^{-1})^{-1}=A$$
        \item[b.] If $A$ and $B$ are $n\times n$ invertible matrices, then so is $AB$, and the inverse of $AB$ is the product of the inverses in reverse order. That is, $$(AB)^{-1}=B^{-1}A^{-1}$$
        \item[c.] If $A$ is an invertible matrix, then so is $A^T$, and the inverse of $A^T$ is the transpose of $A^{-1}$. That is, $$(A^T)^{-1}=(A^{-1})^T$$
    \end{enumerate}
\end{thm}

\begin{thm}
    An $n\times n$ matrix $A$ is invertible iff $A$ is row equivalent to $I_n$, and in this case, any sequence of elementary row ops that reduce $A$ to $I_n$ also transforms $I_n$ into $A^{-1}$.
\end{thm}

\begin{thm} (the invertible matrix theorem).
    Let $A$ be a square $n\times n$ matrix. Then the following statements are equivalent.
    \begin{enumerate}
        \item[a.] $A$ is an invertible matrix.
        \item[b.] $A$ is row equivalent to the $n\times n$ identity matrix.
        \item[c.] $A$ has $n$ pivot positions.
        \item[d.] The equation $A\textbf{x}=\textbf{0}$ has only the trivial solution.
        \item[e.] The columns of $A$ form a linearly independent set.
        \item[f.] The transformation $\textbf{x}\mapsto A\textbf{x}$ is one-to-one.
        \item[g.] The equation $A\textbf{x}=\textbf{b}$ has a unique solution for each $\textbf{b}$ in $\mathbb{R}^n$
        \item[h.] The columns of $A$ span $\mathbb{R}^n$.
        \item[i.] The linear transformation $\textbf{x}\mapsto A\textbf{x}$ maps $\mathbb{R}^n$ onto $\mathbb{R}^n$.
        \item[j.] There is an $n\times n$ matrix $C$ such that $CA=I$.
        \item[k.] There is an $n\times n$ matrix $D$ such that $AD=I$.
        \item[l.] $A^T$ is an invertible matrix.
        \item[m.] The columns of \(A\) form a basis for \(\mathbb{R}^n\).
        \item[n.] \(\text{Col }A=\mathbb{R}^n\).
        \item[o.] \(\text{rank }A=n\).
        \item[p.] \(\text{nullity }A=0\).
        \item[q.] \(\text{Nul }A=\{\textbf{0}\}\).
        \item[r.] The number \(0\) is \textit{not} an eigenvalue of \(A\).
    \end{enumerate}
\end{thm}

\begin{thm}
    Let $T:\mathbb{R}^n\to\mathbb{R}^n$ be a linear transformation and let $A$ be the standard matrix for $T$. Then $T$ is invertible iff $A$ is an invertible matrix. In that case, the linear transformation $S$ given by $S(\textbf{x})=A^{-1}\textbf{x}$ is the unique function satisfying equations
    \begin{align}
        S(T(\textbf{x}))=\textbf{x}\hspace{1cm}\text{for all $x$ in $\mathbb{R}^n$}\\
        T(S(\textbf{x}))=\textbf{x}\hspace{1cm}\text{for all $x$ in $\mathbb{R}^n$}
    \end{align}
\end{thm}
    
\section{Determinants}

\begin{thm}
    The determinant of an \(n\times n\) matrix \(A\) can be computed by a cofactor expansion across any row or down any column. The expansion across the \(i\)-th row is 
    \[
        \det A=a_{i1}C_{i1}+a_{i2}C_{i2}+\cdots+a_{in}C_{in}.
    \] 
    
    The cofactor expansion down the \(j\)-th column is 
    \[
        \det A=a_{1j}C_{1j}+a_{2j}C_{2j}+\cdots+a_{nj}C_{nj}.
    \]

    In either case, \(C_{ij}=(-1)^{i+j}\det A_{ij}\), where \(A_{ij}\) is the \((n-1)\times(n-1)\) matrix that results from deleting the \(i\)-th row and \(j\)-th column of \(A\).
\end{thm}

\begin{thm}
    If \(A\) is a triangular matrix, then \(\det A\) is the product of the entries on the main diagonal.
\end{thm}

\begin{thm} (row operations).
    Let \(A\) be a square matrix.
    \begin{enumerate}
        \item[a.] If a multiple of one row of \(A\) is added to another row to produce matrix \(B\), then \(\det B=\det A\).
        \item[b.] If two rows of \(A\) are interchanged to produce \(B\), then \(\det B=-\det A\).
        \item[c.] If one row of \(A\) is multiplied by \(k\) to produce \(B\), then \(\det B=k\ \det A\).
    \end{enumerate}
\end{thm}

\begin{thm}
    A square matrix \(A\) is invertible iff \(\det A\neq0\).
\end{thm}

\begin{thm}
    If \(A\) is an \(n\times n\) matrix, then \(\det A^T=\det A\).
\end{thm}

\begin{thm} (multiplicative property).
    If \(A\) and \(B\) are \(n\times n\) matrices, then \(\det AB=\left(\det A\right)\left(\det B\right)\).
\end{thm}

\section{Vector Spaces}

\begin{thm}
    If $\textbf{v}_1,\ldots,\textbf{v}_p$ are in a vector space $V$, then Span\{$\textbf{v}_1,\ldots,\textbf{v}_p$\} is a subspace of $V$.
\end{thm}

\begin{thm}
    The null space of an $m\times n$ matrix $A$ is a subspace of $\mathbb{R}^n$.
\end{thm}

\begin{thm}
    The column space of an $m\times n$ matrix $A$ is a subspace of $\mathbb{R}^m$.
\end{thm}

\begin{thm}
    An indexed set \(\textbf{v}_1,\ldots,\textbf{v}_p\) of two or more vectors, with \(\textbf{v}_1\neq0\), is linearly dependent iff some \(\textbf{v}_j\) (with \(j>1\)) is a linear combination of the preceding vectors \(\textbf{v}_1,\ldots,\textbf{v}_{j-1}\).
\end{thm}

\begin{thm} (the spanning set theorem).
    Let \(S=\{\textbf{v}_1,\ldots,\textbf{v}_p\}\) be a set in a vector space \(V\), and let \(H=\spn\{\textbf{v}_1,\ldots,\textbf{v}_p\}\).
    \begin{enumerate}
        \item[a.] If \(\textbf{v}_k\) is a linear combination of the remaining vectors in \(S\), then the set formed by \(S-\{\textbf{v}_k\}\) still spans \(H\).
        \item[b.] If \(H\neq0\), some subset of \(S\) is a basis for \(H\).
    \end{enumerate}
\end{thm}

\begin{thm}
    The pivot columns of a matrix \(A\) form a basis for \(\text{Col }A\).
\end{thm}

\begin{thm}
    If two matrices \(A\) and \(B\) are row equivalent, then their row spaces are the same. If \(B\) is in echelon form, the nonzero rows of \(B\) form a basis for the row space of \(A\) as well as for that of \(B\).
\end{thm}

\begin{thm} (the unique representation theorem).
    Let \(\mathcal{B}=\{\textbf{b}_1,\ldots,\textbf{b}_n\}\) be a basis for a vector space \(V\). Then for each \(\textbf{x}\) in \(V\), there exists a unique set of scalars \(c_1,\ldots,c_n\) such that \[\textbf{x}=c_1\textbf{b}_1+\cdots+c_n\textbf{b}_n\]
\end{thm}

\begin{thm}
    Let \(\mathcal{B}=\{\textbf{b}_1,\ldots,\textbf{b}_n\}\) be a basis for a vector space \(V\). Then the coordinate mapping \(\textbf{x}\mapsto\left[\textbf{x}\right]_{\mathcal{B}}\) is a one-to-one linear transformation from \(V\) onto \(\mathbb{R}^n\).
\end{thm}

\begin{thm}
    If a vector space \(V\) has a basis \(\mathcal{B}=\{\textbf{b}_1,\ldots,\textbf{b}_n\}\), then any set in \(V\) containing more than \(n\) vectors must be linearly dependent.
\end{thm}

\begin{thm}
    If a vector space \(V\) has a basis of \(n\) vectors, then every basis of \(V\) must consist of exactly \(n\) vectors.
\end{thm}

\begin{thm}
    Let \(H\) be a subspace of a finite-dimensional vector space \(V\). Any linearly independent set in \(H\) can be expanded, if necessary, to a basis for \(H\). Also, \(H\) if finite-dimensional and \[\text{dim } H\leq\text{dim }V\]
\end{thm}

\begin{thm} (the basis theorem).
    Let \(V\) be a \(p\)-dimensional vector space, \(p\geq1\). Any linearly independent set of exactly \(p\) elements in \(V\) is automatically a basis for \(V\). Any set of exactly \(p\) elements that spans \(V\) is automatically a basis for \(V\).
\end{thm}

\begin{thm} (the rank theorem).
    The dimensions of the column space and the null space of an \(m\times n\) matrix \(A\) satisfy the equation \[\text{rank }A+\text{nullity }A=\text{number of columns in }A\] 
\end{thm}

\begin{thm}
    Let \(\mathcal{B}=\{\textbf{b}_1,\ldots,\textbf{b}_n\}\) and \(\mathcal{C}=\{\textbf{c}_1,\ldots,\textbf{c}_n\}\) be bases of a vector space \(V\). Then there is a unique \(n\times n\) matrix \(\underset{C\leftarrow B}{P}\) such that \[\left[\textbf{x}_\mathcal{C}\right]=\underset{C\leftarrow B}{P}\left[\textbf{x}_\mathcal{B}\right]\] The columns of \(\underset{C \leftarrow B}{P}\) are the \(\mathcal{C}\)-coordinate vectors of the vectors in the basis \(\mathcal{B}\). That is, \[\underset{C\leftarrow B}{P}=\left[\left[\textbf{b}_1\right]_\mathcal{C}\ \left[\textbf{b}_2\right]_\mathcal{C}\ \cdots \left[\textbf{b}_n\right]_\mathcal{C}\right]\]
\end{thm}

\section{Eigenvalues and Eigenvectors}

\begin{thm}
    The eigenvalues of a triangular matrix are the entries on its main diagonal.
\end{thm}

\begin{thm}
    If \(\textbf{v}_1,\ldots,\textbf{v}_r\) are eigenvectors that correspond to distinct eigenvalues \(\lambda_1,\ldots,\lambda_r\) of an \(n\times n\) matrix \(A\), then the set \(\{\textbf{v}_1,\ldots,\textbf{v}_r\}\) is linearly independent.
\end{thm}

\begin{thm} (properties of determinants).
    Let \(A\) and \(B\) be \(n\times n\) matrices.
    \begin{enumerate}
        \item[a.] \(A\) is invertible iff \(\det{A}\neq0\)
        \item[b.] \(\det AB=(\det A)(\det B)\)
        \item[c.] \(\det A^T=\det A\)
        \item[d.] If \(A\) is triangular, then \(\det A\) is the product of the entries on the main diagonal of \(A\).
        \item[e.] A row replacement operation on \(A\) does not change the determinant. A row interchange changes the sign of the determinant. A row scaling also scales the determinant by the same scalar factor.
    \end{enumerate}
\end{thm}

\begin{thm}
    If \(n\times n\) matrices \(A\) and \(B\) are similar, then they have the same characteristic polynomial and hence the same eigenvalues (with the same multiplicities).
\end{thm}

\begin{rem}
    Two matrices \(A\) and \(B\) are \textit{similar} if there exists an invertible matrix \(P\) such that \(A=PBP^{-1}\).
\end{rem}

\begin{thm} (the diagonalization theorem).
    An \(n\times n\) matrix \(A\) is diagonalizable iff \(A\) has \(n\) linearly independent eigenvectors. In fact, \(A=PDP^{-1}\), with \(D\) a diagonal matrix, iff the columns of \(P\) are the \(n\) linearly independent eigenvectors of \(A\). In this case, the diagonal entries of \(D\) are the eigenvalues of \(A\) that correspond, respectively, to the eigenvectors in \(P\).
\end{thm}

\begin{thm}
    An \(n\times n\) matrix with \(n\) distinct eigenvalues is diagonalizable.
\end{thm}

\begin{thm}
    Let \(A\) be an \(n\times n\) matrix whose distinct eigenvalues are \(\lambda_1,\ldots,\lambda_p\).
    \begin{enumerate}
        \item[a.] For \(1\leq k\leq p\), the dimension of the eigenspace for \(\lambda_k\) is less than or equal to the multiplicity of the eigenvalue \(\lambda_k\).
        \item[b.] The matrix \(A\) is diagonalizable iff the sum of the dimensions of the eigenspaces equals \(n\), and this happens iff \((i)\) the characteristic polynomial factors completely into linear factors and \((ii)\) the dimension of the eigenspace for each \(\lambda_k\) equals the multiplicity of \(\lambda_k\).
        \item[c.] If \(A\) is diagonalizable and \(\mathcal{B}_k\) is a basis for the eigenspace corresponding to \(\lambda_k\) for each \(k\), then the total collection of vectors in \(\mathcal{B}_1,\ldots,\mathcal{B}_p\) forms an eigenvector basis for \(\mathbb{R}^n\).
    \end{enumerate}
\end{thm}

\begin{rem}
    The \emph{trace} of an \(n\times n\) matrix \(A\), denoted by \(\text{tr }A\), is the sum of the diagonal entries of \(A\). If \(\lambda_1,\ldots,\lambda_n\) are the eigenvalues of \(A\), then \(\text{tr }A=\lambda_1+\lambda_2+\cdots+\lambda_n\).
\end{rem}

\begin{rem}
    The determinant of an \(n\times n\) matrix \(A\) is the product of the eigenvalues of \(A\): \[\det A=\lambda_1\lambda_2\cdots\lambda_n\]
\end{rem}

\section{Orthogonality and Least Squares}

\begin{thm} Let \(\textbf{u}\), \(\textbf{v}\), and \(\textbf{w}\) be vectors in \(\mathbb{R}^n\), and let \(c\) be a scalar. Then
    \begin{enumerate}
        \item[a.] \(\textbf{u}\cdot\textbf{v}=\textbf{v}\cdot\textbf{u}\)
        \item[b.] \((\textbf{u}+\textbf{v})\cdot\textbf{w}=\textbf{u}\cdot\textbf{w}+\textbf{v}\cdot\textbf{w}\)
        \item[c.] \((c\textbf{u})\cdot\textbf{v}=c(\textbf{u}\cdot\textbf{v})=\textbf{u}\cdot(c\textbf{v})\)
        \item[d.] \(\textbf{u}\cdot\textbf{u}\geq0\), \hspace{.5cm} and \(\textbf{u}\cdot\textbf{u}=0\iff\textbf{u}=\textbf{0}\)
    \end{enumerate}
\end{thm}

\begin{thm} (The Pythagorean Theorem). 
    Two vectors \(\textbf{u}\) and \(\textbf{v}\) are orthogonal iff \(||\textbf{u}+\textbf{v}||^2=||\textbf{u}||^2+||\textbf{v}||^2\).
\end{thm}

\begin{thm} 
    Let \(A\) be an \(m\times n\) matrix. The orthogonal complement of the row space of \(A\) is the null space of \(A\), and the orthogonal complement of the column space of \(A\) is the null space of \(A^T\): \[(\text{Row }A)^{\perp}=\text{Nul }A\hspace{.5cm}\text{and}\hspace{.5cm}(\text{Col }A)^{\perp}=\text{Nul }A^T\]
\end{thm}

\begin{thm} 
    If \(S=\{\textbf{u}_1,\ldots,\textbf{u}_p\}\) is an orthogonal set of nonzero vectors in \(\mathbb{R}^n\), then \(S\) is linearly independent and hence is a basis for the subspace spanned by \(S\).
\end{thm}

\begin{thm} 
    Let \(\{\textbf{u}_1,\ldots,\textbf{u}_p\}\) be an orthogonal basis for a subspace \(W\) of \(\mathbb{R}^n\). For each \(\textbf{y}\) in \(W\), the weights in the linear combination \[\textbf{y}=c_1\textbf{u}_1+\cdots+c_p\textbf{u}_p\] are given by \[c_j=\frac{\textbf{y}\cdot\textbf{u}_1}{\textbf{u}_1\cdot\textbf{u}_1}\hspace{1cm}(j=1,\ldots,p)\]
\end{thm}

\begin{thm}
    An \(m\times n\) matrix \(U\) has orthonormal columns iff \(U^T U=I\).
\end{thm} 

\begin{thm}
    Let \(U\) be an \(m\times n\) matrix with orthonormal columns, and let \(\textbf{x}\) and \(\textbf{y}\) be in \(\mathbb{R}^n\). Then
    \begin{enumerate}
        \item[a.] \(||U\textbf{x}||=||\textbf{x}||\)
        \item[b.] \((U\textbf{x})\cdot(U\textbf{y})=\textbf{x}\cdot\textbf{y}\)
        \item[c.] \((U\textbf{x})\cdot(U\textbf{y})=0\iff\textbf{x}\cdot\textbf{y}=0\)
    \end{enumerate}
\end{thm}

\begin{thm} (the orthogonal decomposition theorem).
    Let \(W\) be a subspace of \(\mathbb{R}^n\). Then each \(\textbf{y}\) in \(\mathbb{R}^n\) can be written uniquely in the form \[\textbf{y}=\hat{\textbf{y}}+\textbf{z}\] where \(\hat{\textbf{y}}\) is in \(W\) and \(\textbf{z}\) is in \(W^{\perp}\). In fact, if \(\{\textbf{u}_1,\ldots,\textbf{u}_p\}\) is any orthogonal basis of \(W\), then \[\hat{\textbf{y}}=\frac{\textbf{y}\cdot\textbf{u}_1}{\textbf{u}_1\cdot\textbf{u}_1}\textbf{u}_1+\cdots+\frac{\textbf{y}\cdot\textbf{u}_p}{\textbf{u}_p\cdot\textbf{u}_p}\textbf{u}_p\] and \(\textbf{z}=\textbf{y}-\hat{\textbf{y}}\).
\end{thm}

\begin{thm} (the best approximation theorem).
    Let \(W\) be a subspace of \(\mathbb{R}^n\), let \(\textbf{y}\) be any vector in \(\mathbb{R}^n\), and let \(\hat{\textbf{y}}\) be the orthogonal projection of \(\textbf{y}\) onto \(W\). Then \(\hat{\textbf{y}}\) is the closest point in \(W\) to \(\textbf{y}\), in the sense that \[||\textbf{y}-\hat{\textbf{y}}||<||\textbf{y}-\textbf{v}||\] for all \(\textbf{v}\) in \(W\) distinct from \(\hat{\textbf{y}}\).
\end{thm}

\begin{thm}
    If \(\{\textbf{u}_1,\ldots,\textbf{u}_p\}\) is an orthonormal basis for a subspace \(W\) of \(\mathbb{R}^n\), then \[\text{proj}_{W}\textbf{y}=(\textbf{y}\cdot\textbf{u}_1)\textbf{u}_1+(\textbf{y}\cdot\textbf{u}_2)\textbf{u}_2+\cdots+(\textbf{y}\cdot\textbf{u}_p)\textbf{u}_p\] If \(U=\left[\textbf{u}_1\hspace{.25cm}\textbf{u}_2\hspace{.25cm}\cdots\hspace{.25cm}\textbf{u}_p\right]\), then \[\text{proj}_{W}\textbf{y}=U U^T\textbf{y}\] for all \(y\) in \(\mathbb{R}^n\).
\end{thm}

\begin{thm} (Gram-Schmidt).
    Given a basis \(\{\textbf{x}_1,\ldots,\textbf{x}_p\}\) for a nonzero subspace \(W\) of \(\mathbb{R}^n\), define \begin{align*}
    \textbf{v}_1&=\textbf{x}_1\\
    \textbf{v}_2&=\textbf{x}_2-\frac{\textbf{x}_2\cdot\textbf{v}_1}{\textbf{v}_1\cdot\textbf{v}_1}\textbf{v}_1\\
    \textbf{v}_3&=\textbf{x}_3-\frac{\textbf{x}_3\cdot\textbf{v}_1}{\textbf{v}_1\cdot\textbf{v}_1}\textbf{v}_1-\frac{\textbf{x}_3\cdot\textbf{v}_2}{\textbf{v}_2\cdot\textbf{v}_2}\textbf{v}_2\\
    &\vdots\\
    \textbf{v}_p&=\textbf{x}_p-\frac{\textbf{x}_p\cdot\textbf{v}_1}{\textbf{v}_1\cdot\textbf{v}_1}\textbf{v}_1-\frac{\textbf{x}_p\cdot\textbf{v}_2}{\textbf{v}_2\cdot\textbf{v}_2}\textbf{v}_2-\cdots-\frac{\textbf{x}_p\cdot\textbf{v}_{p-1}}{\textbf{v}_{p-1}\cdot\textbf{v}_{p-1}}\textbf{v}_{p-1}
\end{align*}

    Then \(\{\textbf{v}_1,\ldots,\textbf{v}_p\}\) is an orthogonal basis for \(W\). In addition \[\spn\{\textbf{v}_1,\ldots,\textbf{v}_k\}=\spn\{\textbf{x}_1,\ldots,\textbf{x}_k\}\]\hspace{.5cm} for \(1\leq k\leq p\).
\end{thm}

\begin{thm} (QR factorization).
    If \(A\) is an \(m\times n\) matrix with linearly independent columns, then \(A\) can be factored as \(A=QR\), where \(Q\) is an \(m\times n\) matrix whose columns form an orthonormal basis for \(\text{Col}A\) and \(R\) is an \(n\times n\) upper triangular invertible matrix with positive entries on its diagonal.
\end{thm}

\begin{thm}
    The set of least-squares solutions of \(A\textbf{x}=\textbf{b}\) coincides with the nonempty set of solutions of the normal equations \(A^TA\textbf{x}=A^T\textbf{b}\).
\end{thm}

\begin{thm}
    Let \(A\) be an \(m\times n\) matrix. The following statements are logically equivalent:
    \begin{enumerate}
        \item[a.] The equation \(A\textbf{x}=\textbf{b}\) has a unique least-squares solution for each \(\textbf{b}\) in \(\mathbb{R}^m\).
        \item[b.] The columns of \(A\) are linearly independent.
        \item[c.] The matrix \(A^TA\) is invertible.
    \end{enumerate}
When these statements are true, the least-squares solution \(\hat{\textbf{x}}\) is given by \[\hat{\textbf{x}}=(A^TA)^{-1}A^T\textbf{b}\]
\end{thm}

\begin{thm}
    Given an \(m\times n\) matrix \(A\) with linearly independent columns, let \(A=QR\) be a \(QR\) factorization of \(A\). Then, for each \(\textbf{b}\) in \(\mathbb{R}^m\), the equation \(A\textbf{x}=\textbf{b}\) has a unique least-squares solution, given by \[\hat{\textbf{x}}=R^{-1}Q^T\textbf{b}\]
\end{thm}

\section{Symmetric Matrices and Quadratic Forms}

\begin{thm}
    If \(A\) is symmetric, then any two eigenvectors from different eigenspaces are orthogonal.
\end{thm}

\begin{thm}
    An \(n\times n\) matrix \(A\) is orthogonally diagonalizable iff \(A\) is a symmetric matrix.
\end{thm}

\begin{thm} (the pricipal axes theorem).
    Let \(A\) be an \(n\times n\) symmetric matrix. Then there is an orthogonal change of variable, \(\textbf{x}=P\textbf{y}\), that transforms the quadratic form \(\textbf{x}^TA\textbf{x}\) into a quadratic form \(\textbf{y}^TD\textbf{y}\) with no cross-product term.
\end{thm}

\begin{thm} (quadratic forms and eigenvalues).
    Let \(A\) be an \(n\times n\) symmetric matrix. Then a quadratic form \(Q(\textbf{x})=\textbf{x}^TA\textbf{x}\) is
    \begin{enumerate}
        \item[a.] positive definite (\(Q(\textbf{x})>0\) for all \(\textbf{x}\neq0\)) iff the eigenvalues of \(A\) are all positive.
        \item[b.] negative definite (\(Q(\textbf{x})<0\) for all \(\textbf{x}\neq0\)) iff the eigenvalues of \(A\) are all negative.
        \item[c.] indefinite (\(Q(\textbf{x})\) assumes both positive and negative values) iff \(A\) has both positive and negative eigenvalues.
    \end{enumerate}
\end{thm}

\begin{thm}
    Let \(A\) be a symmetric matrix, and define \(m\) and \(M\) as \(m=\min\{\textbf{x}^TA\textbf{x}:|\textbf{x}|=1\}\) and \(M=\max\{\textbf{x}^TA\textbf{x}:|\textbf{x}|=1\}\). Then \(M\) is the greatest eigenvalue \(\lambda_1\) of \(A\) and \(m\) is the least eigenvalue of \(A\). The value of \(\textbf{x}^TA\textbf{x}\) is \(M\) when \(\textbf{x}\) is a unit vector \(\textbf{u}_1\) corresponding to \(M\). The value of \(\textbf{x}^TA\textbf{x}\) is \(m\) when \(\textbf{x}\) is a unit vector \(\textbf{u}_2\) corresponding to \(m\).
\end{thm}

\begin{thm}
    Let \(A\), \(\lambda_1\), and \(\textbf{u}_1\) be as in Theorem 6. Then the maximum value of \(\textbf{x}^TA\textbf{x}\) subject to the constraints \[\textbf{x}^T\textbf{x}=1,\quad\textbf{x}^T\textbf{u}_1=0\] is the second greatest eigenvalue, \(\lambda_2\), and this maximum value is attained when \(\textbf{x}\) is an eigenvector \(\textbf{u}_2\) corresponding to \(\lambda_2\).
\end{thm}

\end{document}
